% !TeX program = pdflatex
\documentclass[aspectratio=169,xcolor={dvipsnames,table}]{beamer}

\input{../../_en_preambulo_beamer}
% Comandos y macros estadísticas compartidas para las presentaciones Beamer en inglés (Metropolis)

% Conjuntos numéricos básicos
\newcommand{\R}{\mathbb{R}}
\newcommand{\N}{\mathbb{N}}
\newcommand{\Q}{\mathbb{Q}}
\newcommand{\Z}{\mathbb{Z}}
\newcommand{\set}[1]{\left\{ #1 \right\}}
\newcommand{\abs}[1]{\left|#1\right|}
\newcommand{\norm}[1]{\left\| #1 \right\|}

% Comandos estadísticos y probabilísticos del curso
\newcommand{\Var}{\operatorname{Var}}
\newcommand{\cov}{\operatorname{Cov}}
\newcommand{\corr}{\rho}
\newcommand{\comb}[2]{\begin{pmatrix} #1 \\ #2 \end{pmatrix}}
\newcommand{\s}{\sigma}
\newcommand{\particion}{\mathfrak{P}}
\newcommand{\card}[1]{\#\left( #1 \right)}
\newcommand{\seq}[3]{\set{#1}_{#2}^{#3}}
\newcommand{\E}{\mathbb{E}}
\newcommand{\Prob}{\mathbb{P}}

% Entornos pedagógicos y cajas en inglés académico natural (soportando tanto nombres en inglés como en español)
\newenvironment{discussion}{%
  \begin{alertblock}{Discussion Question}%
}{%
  \end{alertblock}%
}
\newenvironment{discusion}{%
  \begin{alertblock}{Discussion Question}%
}{%
  \end{alertblock}%
}

\newenvironment{reflection}{%
  \begin{alertblock}{Classroom Reflection}%
}{%
  \end{alertblock}%
}
\newenvironment{reflexion}{%
  \begin{alertblock}{Classroom Reflection}%
}{%
  \end{alertblock}%
}

\newenvironment{paradox}{%
  \begin{alertblock}{Exploratory Paradox}%
}{%
  \end{alertblock}%
}
\newenvironment{paradoja}{%
  \begin{alertblock}{Exploratory Paradox}%
}{%
  \end{alertblock}%
}

\newenvironment{motivation}{%
  \begin{exampleblock}{Motivation: Data Science in Practice}%
}{%
  \end{exampleblock}%
}
\newenvironment{motivacion}{%
  \begin{exampleblock}{Motivation: Data Science in Practice}%
}{%
  \end{exampleblock}%
}

\newenvironment{quickchallenge}{%
  \begin{block}{Quick Team Challenge}%
}{%
  \end{block}%
}
\newenvironment{retoclase}{%
  \begin{block}{Quick Team Challenge}%
}{%
  \end{block}%
}


\title{Simple Random Sampling and the Unbiased Sample Variance}
\subtitle{Section 05.01 --- Statistics, Bessel's Correction, and the Finite Population Correction}
\author[J. Castillo Colmenares]{Juliho Castillo Colmenares}
\institute[Tec de Monterrey]{Tecnologico de Monterrey}
\date{\vspace{-1.5cm}}

\begin{document}

% SLIDE 01: Title Page
\begin{frame}[plain]
	\titlepage
\end{frame}

% SLIDE 02: Roadmap
\begin{frame}{Roadmap --- Chapter 05: Sampling Distributions}
	\vspace{-0.15cm}
	\begin{block}{Curricular Structure of Unit 4}
		\footnotesize
		\begin{enumerate}\itemsep=0.03cm
			\item \textbf{\textcolor{TecRojo}{Section 05.01: Simple Random Sampling, Mean, and Unbiased Sample Variance}} $\leftarrow$ \emph{Current focus}
			\item \textcolor{gray}{Section 05.02: Asymptotic Central Limit Theorem}
			\item \textcolor{gray}{Section 05.03: Chi-Squared Distribution and Sample Variance}
			\item \textcolor{gray}{Section 05.04: Student's $t$-Distribution and Small Samples}
			\item \textcolor{gray}{Section 05.05: Fisher-Snedecor $F$-Distribution}
		\end{enumerate}
	\end{block}
	\vspace{-0.2cm}
	\begin{alertblock}{Learning Objective}
		\scriptsize
		Formalize the concept of a \emph{statistic} as a random variable, master the sampling distribution of the mean ($E(\bar{X})=\mu$, $\Var(\bar{X})=\sigma^2/n$), and prove why the sample variance $S^2$ must be divided by $n-1$ (Bessel's correction) to be an unbiased estimator of $\sigma^2$.
	\end{alertblock}
\end{frame}

% SLIDE 03: Motivation
\begin{frame}{From Distributions to Statistics}
	\footnotesize
	\begin{motivacion}
		So far we have studied individual random variables. In statistical inference, we work with \emph{samples}: collections of observations $X_1, \ldots, X_n$. Any function of the sample --- such as the mean or the sample variance --- is itself a random variable, with its own probability distribution.
	\end{motivacion}
	\vspace{0.15cm}
	\pause
	\begin{itemize}\itemsep=0.1cm
		\item \textbf{Statistic:} a function of the sample that does not depend on unknown parameters ($\bar{X}$, $S^2$, $\hat{p}$). \pause
		\item \textbf{Sampling distribution:} the probability distribution of a statistic. \pause
		\item \textbf{Central question:} how do $\bar{X}$ and $S^2$ behave under repeated sampling?
	\end{itemize}
\end{frame}

% SLIDE 03B: Statistic and Sampling Distribution of the Mean
\begin{frame}{Statistics and the Sampling Distribution of the Mean}
	\vspace{-0.15cm}
	\begin{columns}[T]
		\begin{column}{0.48\textwidth}
			\begin{block}{Mean and Variance of $\bar{X}$}
				\footnotesize
				$E(\bar{X})=\mu$, $\Var(\bar{X})=\sigma^2/n$, $\text{SD}(\bar{X})=\sigma/\sqrt{n}$. $\bar{X}$ is unbiased for $\mu$; its spread shrinks as $n$ grows.
			\end{block}
		\end{column}
		\begin{column}{0.48\textwidth}
			\begin{alertblock}{Consistency}
				\footnotesize
				Since $\Var(\bar{X})\to 0$ as $n\to\infty$ and $E(\bar{X})=\mu$ for all $n$, $\bar{X}$ converges in probability to $\mu$ (Chebyshev): $\bar{X}$ is a \emph{consistent} estimator.
			\end{alertblock}
		\end{column}
	\end{columns}
\end{frame}

% SLIDE 04: Unbiased Sample Variance
\begin{frame}{Unbiased Sample Variance: Bessel's Correction}
	\vspace{-0.15cm}
	\begin{columns}[T]
		\begin{column}{0.48\textwidth}
			\begin{block}{Definition and Theorem}
				\footnotesize
				$S^2 = \frac{1}{n-1}\sum(X_i-\bar{X})^2$, and $E(S^2)=\sigma^2$. The divisor $n-1$ is exactly what makes $S^2$ unbiased.
			\end{block}
		\end{column}
		\begin{column}{0.48\textwidth}
			\begin{alertblock}{Why Not Divide by $n$?}
				\footnotesize
				$\hat{\sigma}^2=\frac{1}{n}\sum(X_i-\bar{X})^2$ has $E(\hat{\sigma}^2)=\frac{n-1}{n}\sigma^2 < \sigma^2$: it systematically underestimates, since $\bar{X}$ is closer to the data than the unknown $\mu$.
			\end{alertblock}
		\end{column}
	\end{columns}
\end{frame}

% SLIDE 05: Proof of Unbiasedness
\begin{frame}{Proof of $E(S^2)=\sigma^2$}
  \footnotesize
  \begin{block}{Key Identity}
    $\sum(X_i-\bar X)^2 = \sum(X_i-\mu)^2 - n(\bar X-\mu)^2$.
  \end{block}
  \pause
  \begin{alertblock}{Taking Expectations}
    $E[\sum(X_i-\bar X)^2] = n\sigma^2 - n(\sigma^2/n) = (n-1)\sigma^2$, hence $E(S^2)=\sigma^2$.
  \end{alertblock}
\end{frame}

% SLIDE 06: Finite Population Correction
\begin{frame}{Sampling Without Replacement: The FPC}
  \footnotesize
  \begin{block}{Finite Population Correction}
    $\Var(\bar X) = \dfrac{\sigma^2}{n}\dfrac{N-n}{N-1}$.
  \end{block}
  \pause
  \begin{alertblock}{Interpretation}
    The factor is below one for finite populations and approaches one as $N\to\infty$; it matters when the sampling fraction is not negligible.
  \end{alertblock}
\end{frame}
% SLIDE 06: Applications
\begin{frame}{Applications in Inference and Data Science}
	\vspace{-0.1cm}
	\begin{itemize}\itemsep=0.1cm
		\item \textbf{Confidence Intervals:} unbiased $S^2$ yields reliable standard errors for $\mu$. \pause
		\item \textbf{Hypothesis Testing:} the $t$-statistic (Section 05.04) depends directly on $S^2$. \pause
		\item \textbf{Auditing and Quality Control:} the FPC adjusts uncertainty when sampling finite batches without replacement. \pause
		\item \textbf{Machine Learning:} unbiased variance in cross-validation and generalization error estimation.
	\end{itemize}
\end{frame}

% SLIDE 07: Python Lab Block 1
\begin{frame}[fragile]{Python Lab: Unbiased Sample Variance}
	\vspace{-0.05cm}
	\scriptsize
	Monte Carlo verification of $E(S^2)=\sigma^2$ vs. the biased estimator, and the shortcut formula (\texttt{05.01\_sample\_statistics.py}):
	{\lstset{basicstyle=\fontsize{4pt}{4.8pt}\selectfont\ttfamily,aboveskip=0.1em,belowskip=0.1em}
	\lstinputlisting[firstline=17, lastline=43]{../../code/05_distribuciones_muestreo/05.01_sample_statistics.py}}
\end{frame}

% SLIDE 08: Python Lab Block 2
\begin{frame}[fragile]{Python Lab: Sampling Distribution of the Mean}
	\vspace{-0.05cm}
	\scriptsize
	Verification of $E(\bar{X})=\mu$ and $\Var(\bar{X})=\sigma^2/n$ across sample sizes:
	{\lstset{basicstyle=\fontsize{4pt}{4.8pt}\selectfont\ttfamily,aboveskip=0.1em,belowskip=0.1em}
	\lstinputlisting[firstline=46, lastline=68]{../../code/05_distribuciones_muestreo/05.01_sample_statistics.py}}
\end{frame}

% SLIDE 09: Python Lab Block 3
\begin{frame}[fragile]{Python Lab: Finite Population Correction}
	\vspace{-0.05cm}
	\scriptsize
	Naive variance vs. the FPC, with empirical verification via sampling without replacement:
	{\lstset{basicstyle=\fontsize{4pt}{4.8pt}\selectfont\ttfamily,aboveskip=0.1em,belowskip=0.1em}
	\lstinputlisting[firstline=71, lastline=102]{../../code/05_distribuciones_muestreo/05.01_sample_statistics.py}}
\end{frame}

% SLIDE 10: Python Lab Terminal Output
\begin{frame}[fragile]{Python Lab: Terminal Output}
	\vspace{-0.1cm}
	\begin{block}{}
		\fontsize{4pt}{5pt}\selectfont
		\begin{verbatim}
=== Block 1: Unbiased Sample Variance (Bessel's Correction) ===
Population: N(50.0, 100.0), n=5, trials=200000
  E(S^2) [unbiased]: 100.1676 (theoretical: 100.0000)
  E(sigma_hat^2) [biased]: 80.1341 (theoretical: 80.0000)
  Worked example {8,12,10,14,6}: mean=10.0000, S^2=10.0000

=== Block 2: Sampling Distribution of the Mean ===
  n=25:  Var(Xbar)=63.98 (theo 64.00) | n=100: 16.04 (theo 16.00) | n=400: 4.00 (theo 4.00)

=== Block 3: Finite Population Correction (FPC) ===
N=200, n=20: naive Var=34.83, FPC factor=0.9045, FPC Var=31.50
Empirical Var(Xbar) [without replacement]: 31.61
		\end{verbatim}
	\end{block}
\end{frame}


% SLIDE 19: Closing and Chapter Perspective
\begin{frame}{Closing Section and Chapter Perspective}
	\vspace{-0.2cm}
	\begin{block}{Synthesis: Statistics and Unbiased Estimation}
		\scriptsize
		Every statistic computed from a sample --- the mean, the variance, the proportion --- is itself a random variable with its own sampling distribution. Formalizing the unbiasedness of $\bar{X}$ and $S^2$ is the first step toward building rigorous confidence intervals and hypothesis tests in the sections ahead.
	\end{block}
	\vspace{0.05cm}
	\pause
	\begin{alertblock}{Key Lessons and Next Step}
		\begin{itemize}\itemsep=0.05cm
			\item \textbf{Sampling distribution of $\bar{X}$:} $E(\bar{X})=\mu$, $\Var(\bar{X})=\sigma^2/n$; unbiased and consistent.
			\item \textbf{Bessel's correction:} $S^2$ divides by $n-1$, not $n$, to be unbiased: $E(S^2)=\sigma^2$.
			\item \textbf{FPC:} sampling without replacement from a finite population multiplies $\Var(\bar{X})$ by $(N-n)/(N-1) < 1$.
			\item \textbf{Next Section:} \textcolor{TecRojo}{Asymptotic Central Limit Theorem}.
		\end{itemize}
	\end{alertblock}
\end{frame}

% SLIDE 20: Modular Perspective
\begin{frame}{Modular Perspective --- The Next Challenge}
  \small
  \begin{block}{What We Have Built}
    Sampling distributions connect descriptive summaries to uncertainty quantification.
  \end{block}
  \pause
  \begin{alertblock}{Next Step}
    The next step is to study asymptotic convergence and the central limit theorem.
  \end{alertblock}
\end{frame}\end{document}
